\documentclass[11pt]{amsart}
\usepackage[margin=1.05in]{geometry}
\usepackage[T1]{fontenc}
\usepackage{lmodern}
\usepackage{amssymb}
\usepackage{booktabs}
\usepackage{enumitem}
\usepackage[hidelinks]{hyperref}
\setlength{\emergencystretch}{2em}
\setlist{leftmargin=2em}
\raggedbottom

\newtheorem{theorem}{Theorem}[section]
\newtheorem{lemma}[theorem]{Lemma}
\newtheorem{proposition}[theorem]{Proposition}
\newtheorem{corollary}[theorem]{Corollary}
\theoremstyle{definition}
\newtheorem{definition}[theorem]{Definition}
\newtheorem{question}[theorem]{Question}
\theoremstyle{remark}
\newtheorem{remark}[theorem]{Remark}

\newcommand{\Z}{\mathbb{Z}}
\newcommand{\C}{\mathcal{C}}

\begin{document}

\title{Multicolour Ramsey numbers of the $3$-uniform loose triangle}

\author{Arhan Canli}
\address{Quantavane, Dubai, United Arab Emirates}
\email{arhan@quantavane.com}

\subjclass[2020]{Primary 05D10; Secondary 05C65, 05C55, 68V15}
\keywords{hypergraph Ramsey numbers, loose cycles, multicolour Ramsey numbers, SAT, DRAT proofs}

\begin{abstract}
Let $R_r(\C)$ be the least $n$ such that every $r$-colouring of the triples of an $n$-set contains a monochromatic
loose triangle $\C$: three triples that pairwise meet in exactly one vertex and span six vertices.
Gy\'arf\'as and Raeisi proved $r+5\le R_r(\C)\le 3r$ for $r\ge3$ and $R_2(\C)=7$, $R_3(\C)=8$, and suspected that
the lower bound $r+5$ is the truth for every $r$; Mubayi and Suk also wrote that the lower bound appears more likely
to be the answer. We show that the suspicion holds for four colours, $R_4(\C)=9$ (a computer proof certified in the
DRAT format), and fails for every $r\ge5$, by an unbounded amount: $R_5(\C)=11$ (an explicit colouring and a
certified cube-and-conquer proof), $R_{12}(\C)\ge19$, and
$R_r(\C)\ge r+6+\lfloor\log_2(r/9)\rfloor$ for every $r\ge9$. The constructions rest on a short lemma about unions of
sunflowers.
\end{abstract}

\maketitle

\section{Introduction}

A \emph{loose triangle} $\C$ (the $3$-uniform loose cycle of length three, $C^3_3$ in \cite{GR}) is the $3$-graph
with vertices $v_1,\dots,v_6$ and edges $\{v_1,v_2,v_3\}$, $\{v_3,v_4,v_5\}$, $\{v_5,v_6,v_1\}$. Equivalently, three
triples form a copy of $\C$ exactly when they pairwise meet in exactly one vertex and no vertex lies in all three.
For $r\ge1$ let $R_r(\C)$ be the least $n$ such that every colouring of the triples of an $n$-set with $r$ colours has
a monochromatic copy of $\C$, and let $f(r)=R_r(\C)-1$, the largest $n$ admitting a \emph{good} colouring (one with no
such copy). A good colouring stays good on every subset of the points.

Gy\'arf\'as and Raeisi \cite[Proposition 13]{GR} proved $r+5\le R_r(\C)\le 3r+1$, the lower bound by colouring
each triple of $\{1,\dots,r+4\}$ by its least element if this is less than $r$ and by $r$ otherwise, and the upper
bound from the Tur\'an number of $\C$ \cite{CK, FF}. Using the uniqueness of the extremal family \cite{CK} they
reduced the upper bound to $3r$ for $r\ge3$; they also proved $R_2(\C)=7$ and $R_3(\C)=8$, and wrote \cite[p.~8]{GR}:
``We proved here that for $r=2,3$ the lower bound is the truth and we suspect that this is true for all $r$.''
Axenovich, Gy\'arf\'as, Liu and Mubayi \cite[Proposition 3]{AGLM} restate the bounds, and Mubayi and Suk
\cite[Section 9]{MS} write that it ``seems to be an interesting open problem to determine which bound above is closer
to the truth. The lower bound appears more likely to be the answer.'' Radziszowski's survey \cite[item 7.2.2(a)]{DS1}
lists $r+5\le R_r(\C)\le 3r$ for $r\ge4$.

For the loose path $\mathcal P$ with three edges the analogous lower bound $r+6$ is the truth for $2\le r\le 10$ (see
\cite[item 7.2.2(b)]{DS1}), and the upper bound has been brought far below the Tur\'an bound \cite{LP, BZ}. For the
loose triangle we find the opposite picture:

\begin{theorem}\label{thm:main}
\begin{enumerate}[label=\textup{(\roman*)}]
\item $R_4(\C)=9$.
\item $R_5(\C)=11$ and $R_{12}(\C)\ge19$.
\item For every $v\ge7$, $R_r(\C)\ge r+v$ whenever $r\ge 9\cdot2^{v-6}-v+1$. In particular
$R_r(\C)\ge r+6+\lfloor\log_2(r/9)\rfloor$ for every $r\ge9$, and $R_r(\C)-r\to\infty$.
\end{enumerate}
\end{theorem}

Part (i) confirms the suspicion of \cite{GR} for $r=4$ and is a computation (Section~\ref{sec:r4}). Parts (ii) and
(iii) refute it for every $r\ge5$; the lower bounds follow from two constructions (Theorems~\ref{thm:odd}
and~\ref{thm:two}), both based on Lemma~\ref{lem:sun}, and the upper bound $R_5(\C)\le11$ is a computation
(Section~\ref{sec:r5}). Thus $R_r(\C)=r+5$ for $r\le4$ and $R_5(\C)=r+6$. Table~\ref{tab:bounds} lists the resulting lower bounds.

\section{Constructions}\label{sec:constr}

\begin{lemma}\label{lem:star}
$f(r+1)\ge f(r)+1$ for every $r\ge1$.
\end{lemma}

\begin{proof}
Add a new point $v$ to a good $r$-colouring of the triples of an $f(r)$-set and give every triple through $v$ a new
colour. Three triples through $v$ share $v$, so the new class has no copy of $\C$.
\end{proof}

A \emph{sunflower} with kernel $e$ (a pair) and petal set $X$ is the set of triples $e\cup\{x\}$, $x\in X$.

\begin{lemma}\label{lem:sun}
Let the point set be partitioned as $P\cup Q$. Let $E_P$ be a set of pairwise disjoint pairs inside $P$ and $E_Q$ a set
of pairwise disjoint pairs inside $Q$, and write $K_P=\bigcup E_P$, $K_Q=\bigcup E_Q$. For each $e\in E_P$ let
$X_e\subseteq (P\cup Q)\setminus K_P$, and for each $e\in E_Q$ let $X_e\subseteq (P\cup Q)\setminus K_Q$. Then
$F=\{e\cup\{x\}: e\in E_P\cup E_Q,\ x\in X_e\}$ contains no copy of $\C$.
\end{lemma}

\begin{proof}
Suppose $t_i=e_i\cup\{x_i\}$ $(i=1,2,3)$ form a copy of $\C$. The kernels are distinct, since two triples with the same
kernel share two points. Two kernels lie on the same side; as the hypotheses are symmetric in $P$ and $Q$ we may
assume $e_1,e_2\in E_P$. They are disjoint, and $x_1\notin e_2$, $x_2\notin e_1$ because the petals avoid $K_P$. Hence
$t_1\cap t_2=\{x_1\}\cap\{x_2\}$, and as $|t_1\cap t_2|=1$ we get $x_1=x_2=:x$. If $e_3\in E_P$ as well, the same
argument gives $x_3=x$, a point common to all three triples. So $e_3\in E_Q$. If $x\in e_3$ or $x_3=x$, again $x$ lies
in all three triples. Otherwise, since $e_1\subseteq P$ and $e_3\subseteq Q$ are disjoint and $x_3\ne x$,
$t_1\cap t_3=(\{x_3\}\cap e_1)\cup(\{x\}\cap e_3)=\{x_3\}\cap e_1$. As $|t_1\cap t_3|=1$, $x_3\in e_1$, and in the
same way $x_3\in e_2$, which is impossible because $e_1\cap e_2=\emptyset$.
\end{proof}

For odd $m$, call a $3$-subset of $\Z_m$ a \emph{progression} if it is $\{c-d,c,c+d\}$ for some $c$ and $d\ne0$ (we say
$c$ is a middle term), and let $U_m$ be the set of $3$-subsets of $\Z_m$ that are not progressions.

\begin{theorem}\label{thm:odd}
Let $m\ge3$ be odd. If the triples in $U_m$ can be coloured with $j$ colours without a monochromatic copy of $\C$, then
$f(m+j)\ge 2m$.
\end{theorem}

\begin{proof}
Take the points $\Z_m\times\{0,1\}$. A triple with two points $(a,s),(b,s)$ in one part and one point in the other
part gets colour $(a+b)/2\in\Z_m$. A progression inside one part gets one of its middle terms (it has one or, when
$3\mid m$, possibly three), and the triples of $U_m$ inside each part get the given $j$-colouring, with $j$ new colours,
the same on both parts. In part $s$ the pairs $\{c-d,c+d\}$, $1\le d\le (m-1)/2$, are disjoint and cover
$\Z_m\setminus\{c\}$, and $\{a,b\}$ is one of them iff $(a+b)/2=c$. So class $c\in\Z_m$ consists of triples $e\cup\{x\}$
with $e$ one of these pairs in a part and $x$ either in the other part or equal to $(c,s)$ in the same part, which is
a family as in Lemma~\ref{lem:sun} with $K_P=(\Z_m\setminus\{c\})\times\{0\}$ and $K_Q=(\Z_m\setminus\{c\})\times\{1\}$.
A class among the $j$ new colours consists of triples inside $P$ and inside $Q$; the three triples of a copy of $\C$
pairwise intersect, so such a copy would lie inside one part, against the hypothesis.
\end{proof}

Every $3$-subset of $\Z_5$ is a progression with a unique middle term (the ten sets $\{c-d,c,c+d\}$, $c\in\Z_5$,
$d\in\{1,2\}$, are distinct), so $U_5=\emptyset$ and Theorem~\ref{thm:odd} gives $f(5)\ge10$: colour a cross triple by
$(a+b)/2$ and a triple $\{a,b,c\}$ inside a part by its middle term, which is $2(a+b+c)$ in $\Z_5$. Each of the five
classes has $24$ triples. For $m=9$ the set $U_9$ has $54$ triples, and a computer search found the following good
$3$-colouring (triples of $\Z_9$ written as digit strings):
\begin{center}\small
\begin{tabular}{ll}
colour 1: & 016 017 026 046 056 067 068 127 137 157 167 178 235 238 245 248 358 458\\
colour 2: & 013 023 034 035 037 038 047 125 128 158 247 256 268 347 457 467 478 568\\
colour 3: & 014 025 028 058 124 134 136 145 146 148 236 257 278 346 356 367 368 578
\end{tabular}
\end{center}
That the three classes contain no copy of $\C$ is checked directly. Theorem~\ref{thm:odd} gives $f(12)\ge18$, that is
$R_{12}(\C)\ge19$, which proves Theorem~\ref{thm:main}(ii).

\begin{theorem}\label{thm:two}
If $m$ is even and $m\le f(j)$, then $f(m-1+j)\ge2m$.
\end{theorem}

\begin{proof}
Take two disjoint $m$-sets $P,Q$ and a $1$-factorization $M_1,\dots,M_{m-1}$ of the complete graph on each of them.
A triple with two points $\{a,b\}$ in one part and one point in the other part gets colour $i$, where $\{a,b\}\in M_i$.
Inside $P$, colour the triples by a good $j$-colouring of an $m$-set with $j$ new colours (it exists as $m\le f(j)$),
and colour the triples inside $Q$ by a copy of the same colouring. Every triple lies inside one part or has exactly two
points in one part, so it gets exactly one colour. In a new colour, a copy of $\C$ would lie inside one part, where the
colouring is good. Cross colour $i$ is the family of Lemma~\ref{lem:sun} with $E_P=M_i$ on $P$, $E_Q=M_i$ on $Q$, and
petal sets $X_e=Q$ for $e\subseteq P$ and $X_e=P$ for $e\subseteq Q$.
\end{proof}

\begin{proof}[Proof of Theorem~\textup{\ref{thm:main}(iii)}]
For $v\ge7$ put $r_v=9\cdot2^{v-6}-v+1$; then $r_7=12$ and $r_{v+1}=2r_v+v-2$. We show $f(r_v)\ge r_v+v-1$ by induction
on $v$. For $v=7$ this is $f(12)\ge18$. Given it for $v$, put $j=r_v$ and $m=r_v+v-1=9\cdot2^{v-6}$, which is even and at
most $f(j)$; Theorem~\ref{thm:two} gives $f(m-1+j)\ge2m$, where $m-1+j=r_{v+1}$ and $2m=r_{v+1}+v$. By
Lemma~\ref{lem:star}, $f(r)-r$ does not decrease, so $R_r(\C)-r=f(r)+1-r\ge v$ for all $r\ge r_v$. For $r\ge18$ let
$v=6+\lfloor\log_2(r/9)\rfloor\ge7$; then $r_v\le9\cdot2^{v-6}\le r$. For $9\le r\le17$ the bound $r+6$ follows from
$f(5)\ge10$ and Lemma~\ref{lem:star}.
\end{proof}

\begin{table}[ht]
\centering
\begin{tabular}{lccccccccc}
\toprule
$r$ & 1--4 & 5--11 & 12--28 & 29--63 & 64--134 & 135--277 & 278--564 & 565--1139 & 1140--2290\\
\midrule
$R_r(\C)-r\ \ge$ & 5 & 6 & 7 & 8 & 9 & 10 & 11 & 12 & 13\\
\bottomrule
\end{tabular}
\caption{Lower bounds for $R_r(\C)-r$ from Lemma~\ref{lem:star}, Theorems~\ref{thm:odd} and~\ref{thm:two} (best
choices for each $r$, computed up to $r=3000$). Previously $5$ for every $r$. For $r\le4$ the value $5$ is exact.}
\label{tab:bounds}
\end{table}

\begin{remark}
Lemma~\ref{lem:sun} uses only that the kernels of a class are pairwise disjoint on each side, so in
Theorem~\ref{thm:two} the factorizations of $P$ and $Q$ may be chosen independently. For $m=7$ the set $U_7$
($14$ triples) has no good $2$-colouring, so Theorem~\ref{thm:odd} gives at best $f(10)\ge14$, which is not new.
\end{remark}

\section{Four colours}\label{sec:r4}

The colouring of Gy\'arf\'as and Raeisi gives $f(4)\ge8$. For the upper bound we show that the triples of a $9$-set
have no good $4$-colouring.

\emph{Encoding.} Number the $84$ triples of $\{0,\dots,8\}$ in lexicographic order, $t=0,\dots,83$, and let the Boolean
variable $x_{t,c}$ with index $4t+c+1$ $(0\le c<4)$ mean ``triple $t$ has colour $c$''. The clauses are: for each $t$,
one clause ``at least one colour'' and six pairwise clauses ``at most one colour''; and for each colour $c$ and each of
the $10080=120\binom96$ copies $\{t_1,t_2,t_3\}$ of $\C$, the clause $\neg x_{t_1,c}\vee\neg x_{t_2,c}\vee\neg x_{t_3,c}$.
This is $336$ variables and $40908$ clauses.

\emph{Certificate without symmetry breaking.} Since all colours play the same role, we may relabel the colours so that
they appear in order of first use along the triples $\{0,1,2\},\{3,4,5\},\{6,7,8\}$: the first has colour $0$, the
second colour at most $1$, and the third colour at most $1$ plus the largest colour used before. Adding these
five clauses gives a formula (SHA-256 \texttt{60ff5150\dots}) that CaDiCaL 3.0.1 refutes; its DRAT proof is verified by
\texttt{drat-trim} (version of 21 April 2024): the solver takes about four minutes, the proof has
$139\,292\,847$ bytes (SHA-256 \texttt{89108a5c573d9ceed874d44e5a383229}\-\texttt{86987f791d8d65e0e6157866a1d284ac}) and
its check takes about two and a half minutes. This certificate does not depend on the lex-leader argument below.

\emph{Certificate with lex-leader symmetry breaking.} The group $G$ generated by permutations of the points and of the
colours acts on assignments by permuting variables and preserves the set of good colourings; if a good colouring
exists, the lexicographically least one $x^*$ in its $G$-orbit satisfies $x^*\le_{\mathrm{lex}}\sigma(x^*)$ for every
$\sigma\in G$ \cite{CGLR}. Adding these constraints for the $36$ point transpositions and the $6$ colour transpositions,
each with the comparator over the positions $p_1<\dots<p_s$ moved by $\sigma$ (in the variable order above) and its own
fresh auxiliary variables $y_1,\dots,y_{s-1}$ ($y_0$ true), namely $\neg y_{i-1}\vee\neg x_{p_i}\vee x_{\sigma(p_i)}$,
$\neg y_{i-1}\vee\neg x_{p_i}\vee\neg x_{\sigma(p_i)}\vee y_i$ and $\neg y_{i-1}\vee x_{p_i}\vee x_{\sigma(p_i)}\vee y_i$,
gives $7350$ variables and $61992$ clauses (SHA-256 \texttt{6a44c615d3677191d08dfcc305a595d2}\-\texttt{59ddc80e00554458bd4087e6};
here every $\sigma$ moves $s=168$ positions). CaDiCaL 3.0.1 refutes it in about three seconds, and its DRAT proof
($3\,562\,285$ bytes, SHA-256 \texttt{93abad83ccacf566e4939b3f6966ea08}\-\texttt{b33fd97fdc2f1c261666e6c5c8c29065})
is verified by \texttt{drat-trim}.

Both formulas are produced identically by two separately written generators (in Python and in C), and the same
encodings with $8$ points are satisfiable; the colourings returned, and the Gy\'arf\'as--Raeisi colouring, pass two
separately written checkers. The generators, formulas, witnesses, checkers, cube files and logs are at
\url{https://github.com/arhancanli/quantavane-loose-triangle-ramsey}.

\section{Five colours}\label{sec:r5}

We show that the triples of an $11$-set have no good $5$-colouring.

\emph{Formula.} The encoding of Section~\ref{sec:r4} with $11$ points and $5$ colours ($165$ triples), with
lex-leader constraints for the $55$ point transpositions and the $10$ colour transpositions, and with two families of
clauses that follow from smaller cases. If a colour class has all its triples through one point $v$, the triples
avoiding $v$ form a $4$-coloured copy of the triples of a $10$-set, which has a monochromatic $\C$ since
$R_4(\C)=9$; so every class has a triple avoiding each point. More generally, for $1\le j\le4$ and every set $S$ of
$11-R_{5-j}(\C)$ points and every set $J$ of $j$ colours, some colour of $J$ appears on a triple avoiding $S$, since
otherwise the triples on the remaining $R_{5-j}(\C)$ points use at most $5-j$ colours ($R_1=6$, $R_2=7$, $R_3=8$,
$R_4=9$). These clause families are invariant under point and colour permutations, so the lex-leader argument still
applies. The formula has $23860$ variables and $355775$ clauses (SHA-256 of the clause part
\texttt{d4ab1fce5ad2f39b0a0198f411926f3b}\-\texttt{ee17f1a53e7b70d9a052e6992bf03204}), and is produced identically by
both generators.

\emph{Cube and conquer.} The look-ahead solver \texttt{march\_cu} splits the formula into $14\,905$ cubes (cutoff depth
$14$). The formula together with the negations of all cubes is refuted by CaDiCaL with a DRAT proof verified by
\texttt{drat-trim}, so every solution would satisfy some cube. Each cube was refuted by CaDiCaL with a time limit of
$300$ seconds and a DRAT proof checked by \texttt{drat-trim}; the $53$ nodes that reached the limit were split again
by \texttt{march\_cu}, with a checked DRAT proof that the formula, the node and the negations of its sub-cubes are
unsatisfiable. In total $18\,129$ leaves were refuted and checked, using about $49$ hours of solver time and $26$ hours
of proof checking. An independent re-check of the generators, of the clause families, of the cover proof and of $200$
random leaves and $5$ split covers (re-solved and re-checked) found no error.

\section{Remarks}

The value $R_5(\C)=11$ shows that the lower bound $r+5$ fails already at five colours; the next open case is
$12\le R_6(\C)\le18$.

\begin{question}
What is the order of growth of $R_r(\C)-r$? In particular, is $R_r(\C)=(1+o(1))r$?
\end{question}

\subsection*{Authorship and the role of AI}
The searches, the encodings, the programs and this draft were produced by AI agents (Claude, by Anthropic) working for
Quantavane under the author's direction; the author takes responsibility for the results. The claims rest on the explicit
colourings, the certified proofs and the arguments written here. No mathematician outside Quantavane has reviewed this draft.

\begin{thebibliography}{9}
\bibitem{AGLM} M. Axenovich, A. Gy\'arf\'as, H. Liu and D. Mubayi, Multicolor Ramsey numbers for triple systems,
\emph{Discrete Math.} \textbf{322} (2014), 69--77.
\bibitem{BZ} T. Bohman and E. Zhu, On multicolor Ramsey numbers of triple system paths of length 3,
arXiv:1907.05236.
\bibitem{CK} R. Cs\'ak\'any and J. Kahn, A homological approach to two problems on finite sets,
\emph{J. Algebraic Combin.} \textbf{9} (1999), 141--149.
\bibitem{CGLR} J. Crawford, M. Ginsberg, E. Luks and A. Roy, Symmetry-breaking predicates for search problems,
in \emph{Proc. KR'96}, 1996, 148--159.
\bibitem{FF} P. Frankl and Z. F\"uredi, Exact solution of some Tur\'an-type problems,
\emph{J. Combin. Theory Ser. A} \textbf{45} (1987), 226--262.
\bibitem{GR} A. Gy\'arf\'as and G. Raeisi, The Ramsey number of loose triangles and quadrangles in hypergraphs,
\emph{Electron. J. Combin.} \textbf{19}(2) (2012), \#P30.
\bibitem{LP} T. \L uczak and J. Polcyn, On the multicolor Ramsey number for 3-paths of length three,
\emph{Electron. J. Combin.} \textbf{24}(1) (2017), \#P1.27.
\bibitem{MS} D. Mubayi and A. Suk, A survey of hypergraph Ramsey problems, arXiv:1707.04229.
\bibitem{DS1} S. P. Radziszowski, Small Ramsey numbers, \emph{Electron. J. Combin.} Dynamic Survey DS1, revision 18
(2026).
\end{thebibliography}

\end{document}
